\documentclass[a4paper,11pt]{article}

\usepackage{amsmath}
\usepackage{url}
\usepackage{hyperref}
\usepackage{enumitem}

\title{Milestone 2: Scientific Contextualization and Tradeoff Analysis}
\author{Christian Percival}
\date{}

\begin{document}

\maketitle

\section*{Purpose}

\begin{itemize}
    \item This milestone places the main paper in the larger real-time systems context.
    \item The goal is not only to summarize related work.
    \item The goal is to compare assumptions, approaches, and tradeoffs.
\end{itemize}

\section*{Main Approach from the Assigned Paper}

\begin{itemize}
    \item The paper combines:
    \begin{itemize}
        \item control performance optimization,
        \item schedulability analysis,
        \item rate adaptation,
        \item Constant Bandwidth Server scheduling,
        \item overrun handling.
    \end{itemize}
    \item The goal is to improve control performance during normal execution.
    \item The system must still guarantee hard deadlines when worst-case execution occurs.
\end{itemize}

\section*{Related Approach 1: WCET-Based Scheduling}

\begin{itemize}
    \item Every task is scheduled using its WCET.
    \item Advantage:
    \begin{itemize}
        \item simple to reason about,
        \item safe if WCET estimates are correct,
        \item strong predictability.
    \end{itemize}
    \item Disadvantage:
    \begin{itemize}
        \item pessimistic,
        \item wastes CPU time when actual execution time is much lower than WCET,
        \item can reduce sampling frequency in control systems.
    \end{itemize}
    \item Relation to topic:
    \begin{itemize}
        \item This is the baseline approach that overrun handling tries to improve.
    \end{itemize}
\end{itemize}

\section*{Related Approach 2: Normal-Time Scheduling Without Protection}

\begin{itemize}
    \item Tasks are scheduled based on normal or average execution time.
    \item Advantage:
    \begin{itemize}
        \item better average CPU utilization,
        \item higher task frequency,
        \item better normal-case control performance.
    \end{itemize}
    \item Disadvantage:
    \begin{itemize}
        \item unsafe when the task exceeds its expected time,
        \item overruns can delay other tasks,
        \item hard deadlines can be missed.
    \end{itemize}
    \item Relation to topic:
    \begin{itemize}
        \item This shows why overrun handling is needed.
    \end{itemize}
\end{itemize}

\section*{Related Approach 3: Constant Bandwidth Server}

\begin{itemize}
    \item CBS assigns a budget $Q_s$ and period $T_s$ to a task.
    \item Bandwidth:
    \[
        U_s = \frac{Q_s}{T_s}
    \]
    \item Advantage:
    \begin{itemize}
        \item provides temporal isolation,
        \item limits how much CPU time a task can consume,
        \item useful for variable execution time tasks.
    \end{itemize}
    \item Disadvantage:
    \begin{itemize}
        \item introduces scheduling overhead,
        \item standard CBS may postpone deadlines too far,
        \item hard-deadline guarantees require careful analysis.
    \end{itemize}
    \item Relation to topic:
    \begin{itemize}
        \item CBS is the main server mechanism used as the basis for overrun management.
    \end{itemize}
\end{itemize}

\section*{Related Approach 4: CBShd}

\begin{itemize}
    \item CBShd is an extension of CBS for hard-deadline tasks.
    \item It considers the remaining computation time.
    \item Advantage:
    \begin{itemize}
        \item handles partial overruns more efficiently,
        \item avoids unnecessary large deadline postponements,
        \item improves responsiveness.
    \end{itemize}
    \item Disadvantage:
    \begin{itemize}
        \item requires estimating remaining computation time,
        \item still introduces runtime overhead,
        \item implementation is more complex than simple WCET scheduling.
    \end{itemize}
\end{itemize}

\section*{Related Approach 5: Resource Sharing Protocols}

\begin{itemize}
    \item Real tasks may share resources such as:
    \begin{itemize}
        \item memory buffers,
        \item sensors,
        \item communication buses,
        \item hardware peripherals.
    \end{itemize}
    \item Shared resources can cause blocking.
    \item The Stack Resource Policy can be used to bound blocking times.
    \item Tradeoff:
    \begin{itemize}
        \item resource protocols improve predictability,
        \item but increase analysis complexity.
    \end{itemize}
\end{itemize}

\section*{Tradeoff: Safety vs Efficiency}

\begin{itemize}
    \item WCET-only scheduling:
    \begin{itemize}
        \item high safety,
        \item low efficiency if WCET is rare.
    \end{itemize}
    \item Normal-time scheduling:
    \begin{itemize}
        \item high efficiency,
        \item unsafe without overrun protection.
    \end{itemize}
    \item CBS/CBShd:
    \begin{itemize}
        \item tries to balance both,
        \item improves efficiency while preserving bounded interference.
    \end{itemize}
\end{itemize}

\section*{Tradeoff: Precision vs Runtime Overhead}

\begin{itemize}
    \item Small server budget:
    \begin{itemize}
        \item more precise overrun handling,
        \item more frequent deadline updates,
        \item higher overhead.
    \end{itemize}
    \item Large server budget:
    \begin{itemize}
        \item lower overhead,
        \item less precise handling of small overruns.
    \end{itemize}
\end{itemize}

\section*{Tradeoff: Offline Guarantees vs Online Adaptation}

\begin{itemize}
    \item Offline analysis:
    \begin{itemize}
        \item checks schedulability before runtime,
        \item provides stronger guarantees.
    \end{itemize}
    \item Online adaptation:
    \begin{itemize}
        \item reacts to actual runtime behavior,
        \item improves flexibility,
        \item harder to guarantee.
    \end{itemize}
    \item The assigned paper combines both:
    \begin{itemize}
        \item offline guarantees using WCET,
        \item online overrun handling through server-based scheduling.
    \end{itemize}
\end{itemize}

\section*{Embedded Systems Consequences}

\begin{itemize}
    \item The approach is relevant for embedded systems with variable workloads.
    \item Good fit:
    \begin{itemize}
        \item robotics,
        \item sensor processing,
        \item visual tracking,
        \item radar tracking,
        \item autonomous systems.
    \end{itemize}
    \item Challenges:
    \begin{itemize}
        \item WCET measurement,
        \item scheduler overhead,
        \item limited CPU resources,
        \item resource sharing,
        \item implementation on Linux or microcontrollers.
    \end{itemize}
\end{itemize}

\section*{Selected References and Why They Are Relevant}

\begin{itemize}
    \item Caccamo, Buttazzo, Sha:
    \begin{itemize}
        \item main paper,
        \item rate adaptation,
        \item overrun handling,
        \item CBShd.
    \end{itemize}
    \item Buttazzo book:
    \begin{itemize}
        \item official background,
        \item real-time scheduling,
        \item overload handling,
        \item resource reservations.
    \end{itemize}
    \item Abeni and Buttazzo:
    \begin{itemize}
        \item original CBS reference.
    \end{itemize}
    \item Liu and Layland:
    \begin{itemize}
        \item foundational EDF and real-time scheduling theory.
    \end{itemize}
    \item Baker:
    \begin{itemize}
        \item Stack Resource Policy and resource sharing.
    \end{itemize}
\end{itemize}

\end{document}